\documentclass[11pt,a4paper]{article}
\usepackage[margin=2.3cm]{geometry}
\usepackage{amsmath,amssymb,booktabs,hyperref,enumitem}
\usepackage[T1]{fontenc}
\usepackage{xcolor}
\hypersetup{colorlinks=true,linkcolor=blue!60!black,urlcolor=blue!60!black}
\newcommand{\Q}{\operatorname{Q}}
\newcommand{\erfc}{\operatorname{erfc}}
\newcommand{\EbNo}{E_b/N_0}
\setlength{\parskip}{4pt}

\title{Theory note S0 --- Why the BER of BPSK in AWGN is $\Q\!\left(\sqrt{2E_b/N_0}\right)$}
\author{Chang Young Ham}
\date{September 2026 \quad (accompanies \texttt{src/s0\_bpsk\_awgn\_ber.py})}

\begin{document}
\maketitle

\noindent\textbf{Purpose.}
S0 is the calibration stage of the repository. S1 and S2 reuse its signal model
($r=s+n$, threshold at zero), so the simulator is first required to reproduce a
result whose exact answer is known. This note derives that answer, maps every
symbol to a line of the code, and derives the statistical test the script uses
to decide whether simulation and theory agree. Mathematics used: Gaussian
random variables, conditional probability, the binomial estimator.

\section{Signal model}

One bit $b\in\{0,1\}$ is transmitted per symbol. Binary phase-shift keying maps it to
\[
  s = \begin{cases} +\sqrt{E_b} & b=1 \\ -\sqrt{E_b} & b=0 \end{cases}
\]
so that every symbol carries energy $s^2 = E_b$ (energy per bit).
The channel adds white Gaussian noise. After the receiver's matched filter the
decision variable is a single real number
\[
  r = s + n, \qquad n \sim \mathcal N\!\left(0,\ \sigma^2\right).
\]

\paragraph{Where $\sigma^2 = N_0/2$ comes from.}
$N_0$ is the \emph{one-sided} noise power spectral density, so the two-sided
density is $N_0/2$. Passing white noise through a filter with impulse response
$h(t)$ gives an output variance
\[
  \sigma^2 = \int_{-\infty}^{\infty}\frac{N_0}{2}\,|H(f)|^2\,df
           = \frac{N_0}{2}\int_{-\infty}^{\infty}|h(t)|^2\,dt
           = \frac{N_0}{2}
\]
for a unit-energy matched filter (the middle step is Parseval's relation).
Using $\sigma^2=N_0$ instead shifts the whole simulated curve by exactly 3\,dB.

\section{The decision rule and the error event}

Both symbols are equally likely, and the noise is symmetric, so the
maximum-likelihood decision is a threshold at zero:
\[
  \hat b = \begin{cases} 1 & r>0 \\ 0 & r\le 0. \end{cases}
\]
Condition on $b=1$ (so $s=+\sqrt{E_b}$). An error occurs when $r\le 0$, i.e.\ when
$n \le -\sqrt{E_b}$:
\[
  P(\text{error}\mid b=1) = P\!\left(n\le-\sqrt{E_b}\right)
  = P\!\left(\frac{n}{\sigma} \le -\frac{\sqrt{E_b}}{\sigma}\right)
  = \Q\!\left(\frac{\sqrt{E_b}}{\sigma}\right),
\]
where $\Q(x) = P(Z > x)$ for $Z\sim\mathcal N(0,1)$ is the Gaussian tail function and we
used the symmetry $P(Z\le -x)=P(Z\ge x)$. By the same argument
$P(\text{error}\mid b=0)$ is identical, so the average bit error rate is
\begin{equation}
  \mathrm{BER} = \Q\!\left(\frac{\sqrt{E_b}}{\sigma}\right)
  = \Q\!\left(\sqrt{\frac{E_b}{N_0/2}}\right)
  = \boxed{\ \Q\!\left(\sqrt{\tfrac{2E_b}{N_0}}\right)\ }.
  \label{eq:ber}
\end{equation}

\paragraph{Numerical check.}
The last point of the S0 sweep is $\EbNo = 9$\,dB $=7.94$. With $E_b=1$,
$\sigma = \sqrt{1/15.89}=0.2509$ and the argument of $\Q$ is $1/0.2509 = 3.986$,
so $\mathrm{BER}=\Q(3.986) = 3.36\times10^{-5}$, the value in
\texttt{results/s0\_ber\_results.csv}. At $\EbNo=10$\,dB, the thermal-noise
condition used in S1, the same calculation gives $\Q(4.47)=3.9\times10^{-6}$.

\section{$\Q$ and $\erfc$}

The code uses \texttt{scipy.special.erfc} because SciPy does not provide $\Q$
directly. The two are related by
\[
  \Q(x) = \tfrac12\,\erfc\!\left(\tfrac{x}{\sqrt2}\right)
  \quad\Longrightarrow\quad
  \mathrm{BER} = \tfrac12\,\erfc\!\left(\sqrt{\tfrac{E_b}{N_0}}\right),
\]
which is \texttt{ber\_bpsk\_theory()} in the script.

\section{Correspondence with the code}

The script fixes $E_b = 1$, so $N_0 = 1/(\EbNo)$ and $\sigma = \sqrt{N_0/2}$.

\begin{center}
\begin{tabular}{@{}ll@{}}
\toprule
Mathematics & Line in \texttt{s0\_bpsk\_awgn\_ber.py} \\
\midrule
$\EbNo$ (linear) from dB & \texttt{ebn0\_lin = db\_to\_linear(ebn0\_db)} \\
$\sigma = \sqrt{1/(2\,\EbNo)}$ & \texttt{sigma = np.sqrt(1.0 / (2.0 * ebn0\_lin))} \\
$b\in\{0,1\}$ & \texttt{bits = rng.integers(0, 2, size=...)} \\
$s = \pm 1$ & \texttt{symbols = 2 * bits - 1} \\
$n\sim\mathcal N(0,\sigma^2)$ & \texttt{noise = rng.normal(0.0, scale=sigma, ...)} \\
$r = s+n$ & \texttt{received = symbols + noise} \\
$\hat b = [r>0]$ & \texttt{decisions = (received > 0).astype(int)} \\
count $\hat b \ne b$ & \texttt{n\_errors += np.sum(decisions != bits)} \\
Eq.~\eqref{eq:ber} & \texttt{0.5 * erfc(np.sqrt(ebn0\_linear))} \\
$z$-score, Eq.~\eqref{eq:z} & \texttt{(n\_err - expected) / np.sqrt(expected)} \\
\bottomrule
\end{tabular}
\end{center}

Note that \texttt{scale} in \texttt{rng.normal} is the standard deviation
$\sigma$, not the variance $\sigma^2$.

\section{Accuracy of the Monte Carlo estimate and the pass criterion}

The simulator observes $K$ errors in $n$ bits and reports $\widehat{\mathrm{BER}} = K/n$.
Each bit is an independent Bernoulli trial with success probability $p=\mathrm{BER}$,
so $K\sim\mathrm{Binomial}(n,p)$ and
\[
  \mathbb E[K] = np, \qquad
  \operatorname{Var}[K] = np(1-p)\approx np .
\]
The \emph{relative} standard deviation of the estimate is therefore
$1/\sqrt{np} = 1/\sqrt{\mathbb E[K]}$.

\paragraph{Stopping rule.}
The script simulates in chunks of $10^6$ bits and stops a point once at least
\texttt{MIN\_ERRORS = 200} errors have been seen. This bounds the relative
spread by $1/\sqrt{200}\approx7\%$ at the \emph{hardest} point: at 9\,dB the
required $n\approx200/p$ is six million bits. The easy points collect far more
errors in their first chunk (79{,}271 at 0\,dB, a spread of $0.36\%$), so the
accuracy differs by a factor of twenty across the sweep.

\paragraph{Pass criterion.}
Because of that difference a single tolerance on the ratio
$\widehat{\mathrm{BER}}/\mathrm{BER}$ would be either too loose at 0\,dB or too
tight at 9\,dB. The script instead standardises each point,
\begin{equation}
  z = \frac{K - np}{\sqrt{np}},
  \label{eq:z}
\end{equation}
which is approximately $\mathcal N(0,1)$ at every point if the simulator is
correct, and requires $|z|<4$. For a correct simulator the probability of
failing at a given point is $P(|Z|>4) = 6\times10^{-5}$, below $10^{-3}$ for the
whole ten-point sweep. The observed values lie between $-2.11$ and $+2.22$,
which is the expected scatter of ten standard normal draws. A modelling error
is detected easily: $\sigma^2=N_0$ instead of $N_0/2$ gives $z>100$ at every
point.

\section{Shape of the curve}

For large $x$ the Gaussian tail satisfies
\[
  \Q(x) \approx \frac{1}{x\sqrt{2\pi}}\,e^{-x^2/2},
\]
so with $x^2 = 2\EbNo$ the BER falls like $e^{-\EbNo}$: exponentially in the
\emph{linear} $\EbNo$. Plotted on a log-$y$ axis against linear $\EbNo$ the
curve would be close to a straight line. Against $\EbNo$ in dB, as in
\texttt{figures/s0\_ber\_curve.png}, equal steps on the $x$-axis are equal
\emph{ratios} of $\EbNo$, so the curve bends downward more and more steeply
(the ``waterfall'' shape). This exponential decay is the benchmark that
partial-band jamming removes in S1, where the decay becomes only
inverse-linear.

\end{document}
